\chapter{YouTube build plan}
\label{app:video}

A production plan for filming this build: series structure, per-episode scripts,
shot list, and the pre-flight checklist that stops a stream dying on air.

\section{The pitch}

\begin{quote}\itshape
``I built a coffee machine that takes contactless payments --- and I wrote every
byte myself, from the SPI bus up. No libraries. Here is how the cryptography in
your bank card actually works.''
\end{quote}

\textbf{Why this works as content:} almost every NFC tutorial online is
\code{import nfc; print(tag.uid)}. Nobody shows the anticollision tree walk, the
mutual authentication, or the reason a debit and a commit are separate commands.
You have a genuinely rare artefact: a from-scratch payment stack that
\emph{works}, with a physical payoff at the end.

\textbf{The hook is the money.} Not ``let us learn OOP'' --- it is \emph{``this
card has \texteuro5 on it, and by the end of this video my code will take
\texteuro1.50 off it.''} Every episode should end with something visibly working.

\section{Series structure}

Two formats, and you should decide up front which you are making.

\begin{center}\small
\begin{tabular}{@{}L{0.18\linewidth}L{0.36\linewidth}L{0.36\linewidth}@{}}
\toprule
 & \thead{Edited series} & \thead{Livestream} \\
\midrule
Length   & 10--25 min per episode & 90--120 min \\
Audience & Discovery, evergreen   & Existing community \\
Prep     & High (script, retakes) & Low --- but rehearse \\
Best for & Episodes 1, 6, 7       & The grind chapters \\
\bottomrule
\end{tabular}
\end{center}

\textbf{Recommendation: livestream the build, then cut a highlights series from
the footage.} You get both, and the edited version benefits from having already
explained everything once.

\subsection{Episode map}

\begin{center}\small
\begin{tabular}{@{}cL{0.34\linewidth}L{0.14\linewidth}L{0.24\linewidth}c@{}}
\toprule
\thead{Ep} & \thead{Title} & \thead{Chapters} & \thead{Payoff on screen} & \thead{Len} \\
\midrule
1  & Talking to an NFC chip with zero libraries & 1--3  & \code{version() = 0x92} & 18 \\
2  & Waking up a contactless card              & 4     & \code{15/15 polls} & 15 \\
3  & Reading a UID --- the anticollision dance  & 5--6  & UID on screen & 20 \\
4  & First words with a DESFire                & 7     & \hex{...91AF} & 12 \\
5  & Writing AES-CMAC from the RFC             & 8     & 4 vectors green & 22 \\
6  & \textbf{Proving I know a secret without sending it} & 9--10 & \code{MUTUAL AUTH OK} & 25 \\
7  & \textbf{Taking money off a real card}     & 11    & \code{5.10 -> 5.00} & 20 \\
8  & Building the kiosk                        & 12, 14 & Tap $\rightarrow$ brew & 18 \\
9  & The phone that tops it up                 & 13    & Balance goes up & 15 \\
10 & What I would do differently               & ---   & Retrospective & 12 \\
\bottomrule
\end{tabular}
\end{center}

Episodes 6 and 7 will carry the series. Put the most effort into episode 6:
mutual authentication is genuinely beautiful and almost nobody explains it well
on video.

\section{Episode scripts}

\subsection{The structure for every episode}

\begin{center}\small
\begin{tabular}{@{}lL{0.62\linewidth}@{}}
\toprule
\thead{Time} & \thead{Beat} \\
\midrule
0:00  & \textbf{COLD OPEN} --- the payoff, out of order. 15 seconds, no talking head. \\
0:15  & \textbf{PROBLEM} --- what does not work yet, concretely \\
1:30  & \textbf{THEORY} --- the idea $+$ the analogy (whiteboard or slide) \\
4:00  & \textbf{BUILD} --- live coding, the hard part \\
12:00 & \textbf{IT FAILS} --- keep this. Always keep this. \\
14:00 & \textbf{FIX} --- the diagnosis, out loud \\
16:00 & \textbf{PAYOFF} --- it works, on hardware \\
17:00 & \textbf{NEXT} --- one sentence teasing the next episode \\
\bottomrule
\end{tabular}
\end{center}

\begin{concept}{Keep the failures}
The single most important production note in this appendix. \textbf{The
debugging \emph{is} the content} --- it is what separates you from a tutorial
that hides the mess.

When your MAC fails and you reason your way to ``the counter increments before
the response'', that is the most valuable ninety seconds in the episode. Nobody
learns from watching code work the first time, and nobody believes it either.
\end{concept}

\subsection{Episode 1 --- Talking to an NFC chip with zero libraries}

\textbf{Cold open:} close-up of the card touching the reader, kiosk screen shows
\texteuro5.10 $\rightarrow$ \texteuro5.00. Cut to black. \emph{``By the end of
this series, every byte of that transaction will be code I wrote. Today: making
the chip say hello.''}

\begin{enumerate}
  \item Show the RC522. \pounds3. Show the datasheet --- the register map page.
        \emph{``This is all the documentation there is, and it is enough.''}
  \item Wire it up on camera. \textbf{Say the 3.3\,V warning out loud} --- it is
        the one mistake that costs a viewer their module.
  \item The problem: \emph{``I could plug this in and print things until it
        works. Two minutes per attempt, and when it fails I will not know if it
        is my code, my wiring, or the chip.''}
  \item \textbf{Analogy: the wall socket.} Draw it.
  \item Build \code{SpiLink}, \code{SpiException}, \code{FakeSpiLink}. Run the
        fake test green --- \textbf{with no hardware plugged in.} Hold up the
        unplugged cable to camera. That visual sells the whole idea.
  \item Build \code{Rc522Driver}. Explain the address encoding on screen, one bit
        at a time.
  \item \textbf{Deliberate failure:} run it and get $-110$. \emph{``That is the
        signed byte thing. Java has no unsigned types.''} Fix with \code{\& 0xFF}.
  \item \textbf{Payoff:} on the Pi. \code{version() = 0x92}.
\end{enumerate}

\textbf{Line to land:} \emph{``That \hex{0x92} means the chip is alive and I am
speaking its language. Nothing between my code and the silicon but two bytes of
SPI.''}

\subsection{Episode 5 --- Writing AES-CMAC from the RFC}

The episode that proves the series is serious.

\textbf{Cold open:} scroll RFC 4493 on screen. \emph{``I am going to implement
this. Not import it --- implement it. And then prove it is correct.''}

\begin{enumerate}
  \item What a MAC is. \textbf{Analogy: the wax seal.} Emphasise: it
        authenticates, it does not hide. Different jobs.
  \item Why CMAC exists --- the last-block problem with naive CBC-MAC, in one
        sentence.
  \item Subkey generation on the whiteboard.
  \item Code \code{leftShift} live. \textbf{Get the carry direction wrong} ---
        deliberately, or genuinely, it is a common slip. All four vectors fail
        with plausible-looking garbage.
  \item \emph{``And this is why we do not test crypto on hardware. Look at that
        output --- it is sixteen bytes of noise. Correct output is also sixteen
        bytes of noise. I cannot tell them apart by looking.''}
        \textbf{This is the thesis of the whole series.} Deliver it slowly.
  \item Fix: walk backwards, carry propagates up.
  \item \textbf{Payoff:} four RFC vectors, green.
\end{enumerate}

\textbf{Line to land:} \emph{``Those four numbers were published in 2006 by
people who have never seen my code. My code reproduces them exactly. That is not
`it seems to work' --- that is proof.''}

\subsection{Episode 6 --- Proving I know a secret without sending it}

The best episode. Take your time.

\textbf{Cold open:} \emph{``This card and my Pi both know a sixteen-byte number.
In the next twenty minutes they are going to prove it to each other --- and if
you record every bit of radio between them, you learn nothing.''}

\begin{enumerate}
  \item Set up the problem properly. \emph{``I cannot send the key. I cannot send
        a hash --- you would replay it. So how?''}
  \item \textbf{Analogy: two spies.} Emphasise it must be \textbf{mutual} ---
        one-way lets an impostor collect the passphrase.
  \item Draw the full dance on the whiteboard. This is the money shot of the
        series; animate it in post if you can.
  \item The testability problem: \emph{``This function calls a random number
        generator. So it does something different every time. How do I write a
        known-answer test for that?''} Pause --- let the viewer feel it.
  \item \textbf{Reveal \code{RandomSource}.} \emph{``Randomness is a dependency
        like any other. Inject it, and I can pin it.''}
  \item Show the test class implementing \emph{both} seams. \emph{``My test is
        the fake card AND the fake dice.''}
  \item \textbf{Payoff, in two stages} --- the ordering matters:
        offline vectors green (\emph{``provably correct with no card in the
        room''}), then the real card: \code{MUTUAL AUTH OK. TI = 8071F582}.
  \item \textbf{Run it twice.} Different TI each time. \emph{``That is not a
        replay of my test data. That is a live handshake with a real chip.''}
        Best twenty seconds in the series.
\end{enumerate}

\subsection{Episode 7 --- Taking money off a real card}

\textbf{Cold open:} the balance dropping. Full screen.

\begin{enumerate}
  \item Three requirements: authenticity, freshness, confidentiality.
        \textbf{Analogy: the cheque} --- amount, signature, cheque number.
  \item The counter, and the ordering rule. Say it three times, three ways. It is
        the bug everyone hits.
  \item Build \code{getValue}. \textbf{Deliberately forget the TI one-liner.}
        Auth succeeds, \code{getValue} fails its MAC. Debug on camera --- a great
        sequence, because the failure is one command away from its cause.
  \item \textbf{The debit/commit split.} Stop and explain properly: \emph{``Debit
        does not move the money. If the card leaves the field before the commit,
        nothing happened. That is not a limitation --- that is the entire safety
        design.''} Then \textbf{physically yank the card mid-transaction} and
        show the balance unchanged. Best demo in the series.
  \item \textbf{Payoff:} \code{510 -> 500 cents}, then tap the phone app and show
        the balance really is 5.00 on an independent device.
\end{enumerate}

\textbf{Line to land:} \emph{``Every layer of that --- the register writes, the
anticollision, the CMAC, the session keys, the handshake, the secure messaging
--- I wrote. It took me [N] weeks and it moves real money.''}

\section{Shot list}

\textbf{Cameras.}
\begin{itemize}
  \item \textbf{A} --- screen capture, 1080p60 (OBS; editor plus a terminal).
  \item \textbf{B} --- overhead on the bench: hands, reader, card. A phone on a
        gooseneck is fine.
  \item \textbf{C} \emph{(optional)} --- talking head for intros and outros.
\end{itemize}

\textbf{Recurring shots you will need repeatedly:}
\begin{itemize}
  \item Card approaching the reader, slow, close. \emph{Get ten takes of this
        early} --- you will reuse it constantly.
  \item Terminal output scrolling to a success line.
  \item Whiteboard or tablet diagrams.
  \item Insert of the datasheet register table.
  \item The kiosk screen, straight on, no glare.
\end{itemize}

\textbf{Audio} matters more than video. A lavalier or a USB condenser. Viewers
forgive 720p; they leave over hiss.

\textbf{Screen legibility:} editor at \textbf{16--18\,pt minimum}, high-contrast
theme, and \textbf{zoom in on the specific lines you are discussing} in post.
Nobody can read your normal font size on a phone.

\section{Livestream pre-flight}

Print this. Run it before every stream.

\textbf{Hardware, 30 minutes before:}
\begin{itemize}[label=$\square$]
  \item Pi booted, on the network, IP noted (\textbf{static lease} --- DHCP will
        move it mid-stream)
  \item \code{ssh} session open and tested
  \item RC522 responds: \code{version() = 0x92}
  \item Test card has a \textbf{known balance} --- set it deliberately with the
        phone app
  \item \textbf{A second card, provisioned identically}, in a drawer
  \item Reader taped down. It \emph{will} get nudged on camera.
\end{itemize}

\textbf{Software:}
\begin{itemize}[label=$\square$]
  \item \code{mvn clean compile} passes
  \item \code{mvn test} all green
  \item Vectors file open in a tab --- you \emph{will} be asked to show them
  \item Git committed to a known-good point; \code{git stash} any half-finished work
\end{itemize}

\textbf{The safety net --- the most important item here:}
\begin{itemize}[label=$\square$]
  \item A branch with the chapter \textbf{already working}, ready to
        \code{git checkout} if you get stuck live
\end{itemize}

\begin{concept}{The ten-minute rule}
That safety net is not cheating. Streaming with no escape hatch is how you spend
forty minutes on a typo in front of an audience.

Give yourself a hard rule: \textbf{stuck for ten minutes on air $\rightarrow$
check out the working branch, walk through it, and move on.} Debug it off-stream
and cover the fix at the start of the next episode. Your audience will respect
the discipline more than the struggle.
\end{concept}

\textbf{Stream ergonomics:}
\begin{itemize}[label=$\square$]
  \item Notifications off, phone silent, second monitor off-camera for chat
  \item Water within reach
  \item Say what you are about to do \emph{before} doing it --- silent typing is
        dead air
\end{itemize}

\section{The honest bits --- say these on camera}

Credibility comes from what you are willing to admit.

\begin{description}[leftmargin=0pt, style=nextline]
  \item[1. The keys are hardcoded and that is not secure.] Say it plainly:
        \emph{``Every card in my fleet has the same key, and that key is in the
        APK in plaintext. Anyone who pulls the APK can mint unlimited balance. For
        a bench project with test cards that is fine --- I am telling you so
        nobody copies this into something real.''} Then explain diversification
        and HSMs briefly. Half your technical audience already knows; the ones who
        do not will learn the most important thing in the series.
  \item[2. Cards can be destroyed.] Warn viewers before the provisioning episode.
        Somebody \emph{will} run it against their office badge.
  \item[3. What you would change.] Episode 10: the FSCI nibble bug that cost you
        an evening, the three-second dead time from a bad wait loop, choosing
        wall-clock deadlines over iteration counts. A retrospective on your own
        mistakes is unusually good content --- and it proves you understand the
        thing rather than having merely typed it.
\end{description}

\section{Titles and thumbnails}

\textbf{Titles that work:}
\begin{itemize}
  \item \emph{I wrote a contactless payment system from scratch (no libraries)}
  \item \emph{How your bank card proves it knows a secret --- without sending it}
  \item \emph{Taking real money off an NFC card with code I wrote myself}
\end{itemize}

\textbf{Titles that do not:} anything starting with ``Learning'', ``Tutorial'',
or ``Part 3''. Lead with the artefact.

\textbf{Thumbnails:} the card mid-tap with the balance visible, or a hex dump
with one line circled. Avoid a face plus arrow if you want the technical
audience --- they read it as low-signal.

\textbf{Description boilerplate:} link this book, the repository, the RFC, and
the NXP datasheet. Timestamp the chapters. Pin a comment with the bill of
materials.

\section{Realistic schedule}

\begin{center}\small
\begin{tabular}{@{}L{0.52\linewidth}L{0.32\linewidth}@{}}
\toprule
\thead{Phase} & \thead{Effort} \\
\midrule
Rebuild solo from this book (no camera) & 3--5 evenings \\
Rehearse each episode's build once      & 1 evening each \\
Record                                  & 1 session each \\
Edit (for the edited series)            & 3--4 h per finished episode \\
\bottomrule
\end{tabular}
\end{center}

\textbf{Do the full rebuild before filming anything.} Not because you will forget
--- because the second time through you will know exactly where the interesting
explanations live, and you will have already hit the bugs you want to keep in.
